\documentclass[10pt,stdletter]{newlfm}
\usepackage{charter}
\usepackage{times}

\widowpenalty=1000
\clubpenalty=1000

\newsavebox{\Luiuc}
\sbox{\Luiuc}{%
	\parbox[b]{1.75in}{%
		\vspace{0.5in}%
		\includegraphics[scale=0.8]
		{hovedlogo_eng.pdf}%
	}%
}%
\makeletterhead{Uiuc}{\Lheader{\usebox{\Luiuc}}}

\newlfmP{headermarginskip=10pt}
\newlfmP{sigsize=10pt}
\newlfmP{dateskipafter=30pt}
\newlfmP{addrfromphone}
\newlfmP{addrfromemail}
\PhrPhone{Phone}
\PhrEmail{Email}

\lthUiuc

\namefrom{\vspace{-1.2cm}Trygve Olav Fossum}
\addrfrom{Department of Marine Technology\\
	Otto Nielsens Veg 10\\
	Trondheim, Norway
}
\phonefrom{(+47) 98 678 309}
\emailfrom{\emph{trygve.o.fossum@ntnu.no}}

\addrto{%
Editor-in-Chief\\
Annals of Applied Statistics\\
}

\greetto{Dear Editor,}
\closeline{Yours faithfully,}

\Lfooter{\vspace{-2cm}Norwegian University of Science and Technology\\
Department of Marine Technology, Applied Underwater Robotics Laboratory\\
NO-7491, Trondheim, Norway\\ 
}

\Rheader{1 of 1}

\begin{document}
\begin{newlfm}

We thank the editors and reviewers for providing feedback on our manuscript \textbf{AOAS1910-027}. The feedback have been insightful and constructive. Our response is attached along with the updated manuscript itself. We have carefully considered all the feedback in this revised version of the manuscript, which we believe is a significantly improved version of our work. 

After extensive discussion we collectively felt the need to not only recruit an additional co-author, but also to change the title to 'Learning Excursion Sets of Vector-valued Gaussian Random Fields for Autonomous Ocean Sampling', which we believe better reflects the current manuscript and the valuable reviewer and editorial feedback. 

We look forward to receiving any further feedback at the earliest.
\vspace{1em}

% combining new methods and tools to achieve increased sampling efficiency.
%Towards this end, the work introduces informative descriptions of spatial characteristics of ocean phenomena by define a design criterion based on reducing uncertainty in the joint excursions of random fields, and derive analytical expressions for the expected Bernoulli variance reduction in the excursion set of spatial temperature and salinity variables. 
%Resolving this type of sensing problem require theory and assets that can capitalize on the 
%Combining oceanographic sampling, combining compact models and autonomous robotic platforms for upper water-column observation. The work is novel in multiple ways. First, it uses a combination of unsupervised machine learning and hierarchical clustering methods to distill a compact ocean model based on remote sensing images. The method is capable of characterizing important dominant sea conditions using these images by grouping them into distinct classes. The classification is then verified using preceding wind data to ensure that the model is physically correct in oceanographical sense. Second, the work is supplemented with theory for doing forecast and predictions, as well as evaluation of model error. The combined study demonstrates a novel sampling perspective, combining new methods and tools to achieve increased sampling efficiency.\par \vspace{1em}
%We believe such a work to be representative of new methods and approaches which are necessary and critical to understand our changing oceans.

\end{newlfm}
\end{document}
